\clearpage
\fancyhead[L]{\small ANZEIGEN.}
\SourceLeaf{100}
\pdfbookmark[0]{Anzeigen}{anzeigen}
\begin{center}
{\Large\bfseries ANZEIGEN.}\par
\vspace{1.8em}
\rule{.88\textwidth}{.3pt}\par
Göttingische gelehrte Anzeigen.\quad 26. Februar 1821.\par
\rule{.88\textwidth}{.3pt}
\end{center}
\pdfbookmark[1]{26. Februar 1821}{anzeige-1821-02-26}

Am 15. Februar wurde der Königlichen Gesellschaft von Herrn Hofrat \textsc{Gauss} eine Vorlesung übergeben mit dem Titel
\begin{center}
\textit{Theoria Combinationis observationum erroribus minimis obnoxiae, pars prior,}
\end{center}
\noindent die eine der wichtigsten Anwendungen der Wahrscheinlichkeitsrechnung zum Gegenstand hat. Alle Beobachtungen, die sich auf die Bestimmung von Größen in der sinnlich wahrnehmbaren Welt beziehen, können niemals \textit{absolute} Genauigkeit besitzen, wie genau sie auch angestellt und wie vortrefflich die verwendeten Instrumente auch sein mögen; sie bleiben stets Näherungen, die größeren oder kleineren Fehlern unterliegen. Hier sind nicht solche Fehler gemeint, deren Quellen genau bekannt sind und deren Größe sich bei bestimmten Beobachtungen jedes Mal berechnen lässt. Denn solche Fehler können und sollen von den beobachteten Größen abgezogen werden; das ist dasselbe, als wären sie gar nicht vorhanden. Ganz anders verhält es sich mit den als zufällig anzusehenden Fehlern, die aus der begrenzten Schärfe der Sinne, aus verschiedenen unvermeidlichen und keiner Regel folgenden Unvollkommenheiten der Instrumente sowie aus äußeren Störungen herrühren, die unregelmäßig (jedenfalls für uns) wirken – etwa aus dem Wallen der Atmosphäre beim Sehen oder aus dem Mangel an absoluter Festigkeit beim Aufstellen der Instrumente. Diese zufälligen Fehler, die
\SourceLeaf{101}
sich nicht der Rechnung unterwerfen lassen, können nicht \textit{beseitigt} werden; durch sorgfältige Aufmerksamkeit und durch Vervielfachung der Beobachtungen kann der Beobachter sie nur \textit{vermindern}. Nachdem der Beobachter das Seine getan hat, ist es jedoch Sache des Geometers, die Unsicherheit der Beobachtungen und der daraus durch Rechnung abgeleiteten Größen nach strengen mathematischen Grundsätzen zu beurteilen. Vor allem muss er dort, wo sich mit den Beobachtungen zusammenhängende Größen durch verschiedene Kombinationen aus ihnen ableiten lassen, diejenige Art vorschreiben, bei der die geringstmögliche Unsicherheit zu befürchten ist.

Obwohl zufällige Fehler als solche keinem Gesetz folgen und ohne Ordnung auftreten – in einer Beobachtung größer, in einer anderen kleiner –, steht doch fest: Bei einer bestimmten Beobachtungsart und als fest angesehenen Eigenheiten des Beobachters und seiner Instrumente liegen die aus jeder einzelnen Fehlerquelle hervorgehenden Fehler nicht nur innerhalb gewisser Grenzen; auch jeder mögliche Fehler zwischen diesen Grenzen hat seine bestimmte relative Wahrscheinlichkeit. Je nach ihrer Größe sind solche Fehler also häufiger oder seltener zu erwarten. Wer die Beschaffenheit einer solchen Fehlerquelle genau und vollständig verstünde, könnte diese Grenzen und den Zusammenhang zwischen der Wahrscheinlichkeit der einzelnen Fehler und ihrer Größe bestimmen, ähnlich wie sich bei einem Glücksspiel, sobald seine Regeln bekannt sind, die Grenzen möglicher Gewinne und Verluste und deren relative Wahrscheinlichkeiten berechnen lassen. Dasselbe gilt für den Gesamtfehler, der aus dem Zusammenwirken der einzelnen Fehlerquellen hervorgeht. Diese Begriffe beschränken sich nicht auf unmittelbare Beobachtungen, sondern gelten auch für mittelbare Größenbestimmungen, die aus Beobachtungen abgeleitet werden. In der Praxis werden uns allerdings fast immer die Mittel fehlen, das Gesetz der Fehlerwahrscheinlichkeiten \textit{a priori} anzugeben.

Wie wir im Allgemeinen die Unzulässigkeit einer bestimmten Beobachtungsart beurteilen wollen, hängt teilweise von unserem Ermessen ab. Wir können allein die Größe der äußersten möglichen Fehler zum Maßstab nehmen oder zusätzlich die größere oder geringere Wahrscheinlichkeit der einzelnen möglichen Fehler berücksichtigen. Letzteres scheint angemessener. Diese Berücksichtigung lässt sich jedoch auf vielfache Weise
\SourceLeaf{102}
vornehmen. Bisher haben die Rechner als Maßstab den sogenannten wahrscheinlichen (nicht \textit{wahrscheinlichsten}) Fehler verwendet: Es ist derjenige Fehler, jenseits dessen alle möglichen Fehler zusammengenommen gerade ebenso viel Wahrscheinlichkeit haben wie alle diesseits davon. \textit{Viel vorteilhafter} wird es jedoch sein, zu diesem Zweck anstelle des wahrscheinlichen Fehlers den \textit{mittleren} zu verwenden, sofern dieser an sich noch schwankende Begriff richtig bestimmt wird. Man ordne jedem Fehler ein von seiner Größe abhängiges Moment zu, multipliziere das Moment jedes möglichen Fehlers mit dessen Wahrscheinlichkeit und addiere die Produkte. Als mittlerer Fehler ist dann derjenige anzusehen, dessen Moment diesem Aggregat gleich ist. Welche Funktion der Fehlergröße wir als Moment wählen, bleibt wiederum unserem \textit{Ermessen} überlassen, sofern die Funktion stets positiv ist und für größere Fehler größere Werte annimmt als für kleinere. Der Verfasser hat die einfachste solche Funktion gewählt, nämlich das Quadrat. Diese Wahl ist außerdem mit mehreren höchst wesentlichen Vorteilen verbunden, die keine andere bietet. Sonst könnte auch jede andere Potenz mit geradem Exponenten verwendet werden; je größer der Exponent gewählt wird, desto näher kommt man dem Prinzip, bei dem allein die äußersten Fehler als Maßstab der Genauigkeit dienen. Gegen die Art, wie ein großer Geometer den Begriff des mittleren Fehlers gefasst hat – indem er die Momente den Fehlern gleichsetzt, wenn diese positiv sind, und bei negativen Fehlern die entgegengesetzten Größen verwendet –, lässt sich einwenden, dass dies der mathematischen Stetigkeit zuwiderläuft, ebenso willkürlich gewählt ist wie jede andere Art und weit weniger einfache und befriedigende Ergebnisse liefert. Auch scheint es an sich natürlicher, das Moment der Fehler stärker wachsen zu lassen als die Fehler selbst, denn gewiss lässt man sich lieber den einfachen Fehler zweimal als den doppelten einmal gefallen.

Diese Erläuterungen mussten vorausgeschickt werden, wenn hier auch nur etwas vom Inhalt der Untersuchung angeführt werden sollte, deren ersten Teil die gegenwärtige Abhandlung bildet.

Wenn die Größen, deren Werte durch Beobachtungen gefunden werden, auf bekannte Weise mit einer gleichen Anzahl unbekannter Größen zusammenhängen, lassen sich die Werte der unbekannten Größen\SourceLeaf{103} im Allgemeinen durch Rechnung aus den Beobachtungen ableiten. Natürlich sind diese Werte nur näherungsweise richtig, soweit es auch die Beobachtungen waren. Die Wahrscheinlichkeitsrechnung hat hier lediglich die Unsicherheit dieser Bestimmungen zu beurteilen und setzt dabei die Unsicherheit der Beobachtungen voraus. Gibt es mehr unbekannte Größen als Beobachtungen, lassen sich jene aus diesen überhaupt nicht bestimmen. Gibt es dagegen weniger unbekannte Größen als Beobachtungen, so ist die Aufgabe überbestimmt: Es gibt dann unendlich viele Kombinationen, mit denen sich die unbekannten Größen aus den Beobachtungen ableiten lassen. Wären die Beobachtungen absolut genau, müssten all diese Kombinationen zwar zu denselben Ergebnissen führen; unter den gegebenen Umständen ergeben sie jedoch Ergebnisse, die mehr oder weniger voneinander abweichen. Aus dieser unendlichen Mannigfaltigkeit von Kombinationen die zweckmäßigste auszuwählen – also diejenige, bei der die Unsicherheit der Ergebnisse so klein wie möglich wird –, ist unstreitig eine der wichtigsten Aufgaben bei der Anwendung der Mathematik auf die Naturwissenschaften.

Der Verfasser der gegenwärtigen Abhandlung, der diese Aufgabe 1797 erstmals nach den Grundsätzen der Wahrscheinlichkeitsrechnung untersuchte, erkannte bald, dass sich die \textit{wahrscheinlichsten} Werte der unbekannten Größen nicht ermitteln lassen, wenn die Funktion, welche die Fehlerwahrscheinlichkeit darstellt, unbekannt ist. Solange sie unbekannt ist, bleibt nichts anderes, als eine solche Funktion hypothetisch anzunehmen. Am natürlichsten erschien es ihm, zunächst den umgekehrten Weg einzuschlagen und die Funktion zu suchen, die vorausgesetzt werden muss, damit sich daraus für den einfachsten aller Fälle eine allgemein anerkannte gute Regel ergibt: nämlich, dass das arithmetische Mittel aus mehreren Werten, die durch gleich zuverlässige Beobachtungen für ein und dieselbe unbekannte Größe gefunden wurden, als der wahrscheinlichste gelten muss. Daraus ergab sich, dass die Wahrscheinlichkeit eines Fehlers \(x\) proportional zu einer Exponentialgröße der Form \(e^{-hhxx}\) anzunehmen ist und dass dann gerade die Methode allgemein notwendig wird, zu der er bereits einige Jahre zuvor durch andere Überlegungen gelangt war. Diese Methode, die er später, besonders seit 1801, bei allerlei astronomischen Rechnungen fast täglich anwenden konnte und zu der inzwischen auch \textsc{Legendre} gelangt war,
\SourceLeaf{104}
ist heute unter dem Namen Methode der kleinsten Quadrate allgemein gebräuchlich. Ihre Begründung durch die Wahrscheinlichkeitsrechnung sowie die Bestimmung der Genauigkeit der Ergebnisse selbst und weitere damit zusammenhängende Untersuchungen sind ausführlich in der \textit{Theoria Motus Corporum Coelestium} entwickelt.

Der Marquis \textsc{Delaplace}, der diesen Gegenstand später von einem neuen Gesichtspunkt aus betrachtete – indem er nicht die wahrscheinlichsten Werte der unbekannten Größen suchte, sondern die zweckmäßigste Kombination der Beobachtungen –, gelangte zu dem bemerkenswerten Ergebnis, dass die Methode der kleinsten Quadrate die zweckmäßigste Kombination ist, wenn man die Anzahl der Beobachtungen als unendlich groß betrachtet, und zwar unabhängig von der Funktion, welche die Fehlerwahrscheinlichkeit ausdrückt.

Daran zeigt sich, dass beide Begründungsarten noch etwas zu wünschen übrig lassen. Die erste hängt ganz von der hypothetischen Form der Fehlerwahrscheinlichkeit ab. Verwirft man diese Form, so sind die mit der Methode der kleinsten Quadrate gefundenen Werte der unbekannten Größen tatsächlich nicht mehr die wahrscheinlichsten – ebenso wenig wie die arithmetischen Mittel in dem zuvor angeführten einfachsten Fall. Die zweite Begründungsart lässt uns völlig im Dunkeln darüber, was bei einer mäßigen Zahl von Beobachtungen zu tun ist. Dann hat die Methode der kleinsten Quadrate nicht länger den Rang eines von der Wahrscheinlichkeitsrechnung gebotenen Gesetzes, sondern empfiehlt sich nur durch die Einfachheit der damit verbundenen Rechenoperationen.

Der Verfasser hat diese Untersuchung in der gegenwärtigen Abhandlung erneut aufgenommen. Er ging von einem ähnlichen Gesichtspunkt wie \textsc{Delaplace} aus, bestimmt den zu befürchtenden mittleren Fehler jedoch auf andere Weise, die ihm schon an sich natürlicher erscheint. Er hofft, dass die Freunde der Mathematik mit Vergnügen sehen werden, wie die Methode der kleinsten Quadrate auf der hier gegebenen neuen Grundlage allgemein als die zweckmäßigste Kombination der Beobachtungen erscheint: nicht näherungsweise, sondern mit mathematischer Genauigkeit, unabhängig davon, welche Funktion die Fehlerwahrscheinlichkeit beschreibt und ob die Zahl der Beobachtungen groß oder klein ist.


Mit dem Hauptgegenstand hängt eine Vielzahl weiterer bemerkenswerter Untersuch\-\SourceLeaf{105}ungen eng zusammen. Ihr Umfang zwang den Verfasser jedoch, die Entwicklung der meisten einer künftigen zweiten Vorlesung vorzubehalten. Von denjenigen, die bereits im gegenwärtigen ersten Teil vorkommen, sei es uns erlaubt, hier nur ein Ergebnis anzuführen. Ist die Funktion unbekannt, die die relative Wahrscheinlichkeit jedes einzelnen Fehlers ausdrückt, lässt sich natürlich auch nicht bestimmen, wie wahrscheinlich es ist, dass der Fehler zwischen gegebenen Grenzen liegt. Wenn aber größere Fehler stets eine geringere (zumindest keine größere) Wahrscheinlichkeit haben als kleinere, muss die Wahrscheinlichkeit, dass der Fehler zwischen den Grenzen \(-x\) und \(+x\) liegt, dennoch notwendig größer (zumindest nicht kleiner) sein als \(\frac{x}{m}\surd\frac13\), sofern \(x\) kleiner ist als \(m\surd\frac34\), und nicht kleiner als \(1-\frac{4mm}{9xx}\), sofern \(x\) größer ist als \(m\surd\frac34\). Dabei bezeichnet \(m\) den bei den Beobachtungen zu befürchtenden mittleren Fehler. Für \(x=m\surd\frac34\) stimmen, wie man sieht, beide Ausdrücke überein.
\par\bigskip\centerline{\rule{.88\textwidth}{.3pt}}




